% preamble.tex —— 全书宏包与自定义命令
% 引擎：XeLaTeX；中文：ctex（Windows 字体集）；参考文献：biber/biblatex

\documentclass[11pt,a4paper,openany]{book}
\usepackage[UTF8,fontset=windows]{ctex}
\usepackage{amsmath,amssymb,amsthm}
\usepackage[top=2.6cm,bottom=2.6cm,left=2.6cm,right=2.4cm]{geometry}
\usepackage{xcolor}
\usepackage{graphicx}
\usepackage{booktabs}
\usepackage{longtable}
\usepackage{array}
\usepackage{multirow}
\usepackage{tikz}
\usetikzlibrary{arrows.meta,positioning,calc,shapes.geometric}
\usepackage{pgfplots}
\pgfplotsset{compat=1.18}
\usepackage{listings}
\usepackage[most]{tcolorbox}
\usepackage{enumitem}
\usepackage[hidelinks,bookmarksnumbered]{hyperref}
\usepackage[backend=biber,style=authoryear,sorting=nyt,maxbibnames=4]{biblatex}
% Load after biblatex so its URL penalties allow full commit identifiers to wrap.
\usepackage{xurl}
\addbibresource{bib/references.bib}
\setlength{\bibitemsep}{-0.8pt}
% Parenthesize author-year citations so they remain distinct from Chinese prose.
\let\cite\parencite

% Reserve room for section numbers such as 16.10 in the table of contents.
\makeatletter
\renewcommand*\l@section{\@dottedtocline{1}{1.5em}{3.5em}}
\makeatother

% ---------- 草稿开关：定稿时改为 \draftfalse ----------
\newif\ifdraft
\draftfalse
\newcommand{\TODO}[1]{\ifdraft\textcolor{red}{\textbf{[待写：#1]}}\fi}
\newcommand{\CITECHECK}[1]{\ifdraft\textcolor{orange}{\textbf{[待核来源：#1]}}\fi}

% ---------- 数字身份：cited / calc / measured ----------
% 第一个参数一律用 ASCII 身份（cited / calc / measured），与 tools/check_consistency.py
% 的白名单一致；显示的中文标签由查表给出，避免中文身份与检查规则冲突。
\makeatletter
\@namedef{datalabel@cited}{引用}
\@namedef{datalabel@calc}{计算}
\@namedef{datalabel@measured}{实测}
\newcommand{\datalabel}[1]{\@ifundefined{datalabel@#1}{#1}{\@nameuse{datalabel@#1}}}
\makeatother
\newcommand{\data}[2]{#2\,\raisebox{0.35ex}{\tiny\textcolor{gray!70!black}{[\datalabel{#1}]}}}
\newcommand{\cited}[1]{\data{cited}{#1}}
\newcommand{\calcnum}[1]{\data{calc}{#1}}
\newcommand{\measured}[1]{\data{measured}{#1}}

% ---------- 章首四件套 ----------
\newcommand{\chaptergoal}[1]{\paragraph{本章问题}#1}
\newcommand{\chapterprops}[1]{\paragraph{关联研究问题}#1}
\newcommand{\chapterfig}[1]{\paragraph{核心图}#1}
\newcommand{\chapterpitfall}[1]{\paragraph{本章反例}#1}

% ---------- 六段式 ----------
\newtcolorbox{failurebox}{
  colback=orange!5!white, colframe=orange!60!black,
  title={具体失败}, fonttitle=\bfseries, boxrule=0.4pt, left=6pt, right=6pt, top=4pt, bottom=4pt
}
\newtcolorbox{principlebox}{
  colback=blue!5!white, colframe=blue!55!black,
  title={选择原则}, fonttitle=\bfseries, boxrule=0.4pt, left=6pt, right=6pt, top=4pt, bottom=4pt
}
\newtcolorbox{pitfallbox}{
  colback=red!4!white, colframe=red!55!black,
  title={反例与失效}, fonttitle=\bfseries, boxrule=0.4pt, left=6pt, right=6pt, top=4pt, bottom=4pt
}
\newcommand{\designquestion}[1]{\subsection*{设计问题}\noindent#1}
\newcommand{\falsifier}[1]{\paragraph{检验条件}#1}
\newcommand{\openquestion}[1]{\paragraph{开放问题}#1}

% ---------- 命题环境 ----------
\newtheorem{proposition}{研究问题}[chapter]
\theoremstyle{definition}
\newtheorem{definition}{定义}[chapter]
\newtheorem{example}{例}[chapter]

\newcommand{\sixquestions}{%
  \noindent\textit{读者方法卡：请用具体工件与来源填写以下六问。}\par
  \begin{enumerate}[label=Q\arabic*.]
    \item 改什么：\dotfill
    \item 根据什么改：\dotfill
    \item 怎样改：\dotfill
    \item 什么时候改：\dotfill
    \item 保留什么：\dotfill
    \item 怎样检验：\dotfill
  \end{enumerate}
}

% ---------- 代码样式（listings，不依赖 shell-escape） ----------
\lstset{
  basicstyle=\ttfamily\small,
  frame=single,
  framerule=0.3pt,
  backgroundcolor=\color{gray!6!white},
  keywordstyle=\bfseries,
  commentstyle=\itshape,
  numbers=left,
  numberstyle=\tiny,
  breaklines=true,
  tabsize=2,
  columns=fullflexible,
  keepspaces=true,
  extendedchars=true,
  literate={-}{-}1
}
\lstdefinelanguage{json}{
  basicstyle=\ttfamily\small,
  keywords={true,false,null},
  morestring=[b]"
}

% ---------- 表格：结果表的统一格式 ----------
\newcommand{\resulttablenote}{%
  \begin{footnotesize}
  表中数字身份：[引用] 原论文报告；[计算] 本书计算脚本；[实测] 本书 \texttt{experiments/} 实测。
  \end{footnotesize}
}

\hypersetup{
  pdftitle={深入理解自进化 Agent：经验如何变成可检验的能力变化},
  pdfauthor={shawn},
  pdfsubject={Agent 自进化},
  citecolor=blue!60!black,
  linkcolor=blue!60!black,
  urlcolor=blue!60!black
}
